% relatedwork.tex -- Related Work

The classical Hopfield network \cite{hopfieldNeuralNetworksPhysical1982}
stores binary patterns as local minima of a quadratic energy. Retrieval
proceeds by energy descent, but the storage capacity scales linearly with the
state dimension \cite{amitStoringInfiniteNumbers1985}. Dense associative
memories \cite{krotovDenseAssociativeMemory2016,krotovLargeAssociativeMemory2021}
replace the quadratic interaction with higher-order functions. Modern Hopfield
networks \cite{ramsauerHopfieldNetworksAll2021} use a log-sum-exp energy with
exponential storage capacity and connect its retrieval map to dot-product
attention. Work on random-pattern ensembles and temperature-dependent behavior
has characterized retrieval, condensation, and capacity boundaries
\cite{lucibelloExponentialCapacityDense2024,koulischerExploringTemperatureDependent2023}.
The older Boltzmann-machine literature
\cite{ackleyLearningAlgorithmBoltzmann1985} supplies the complementary idea of
sampling from an associative-memory energy, although it uses Gibbs dynamics on
binary states rather than Langevin dynamics in a continuous state space. The
Gaussian-mixture interpretation developed here is also related to established
probabilistic views of attention. Tsai et al.
\cite{tsaiTransformerKernel2019} describe attention as a kernel smoother,
Nguyen et al.
\cite{nguyenProbabilisticKeys2022} derive attention weights as posterior
responsibilities under Gaussian keys, and Ota and Karakida
\cite{otaAttentionBoltzmannMachine2023} study a Boltzmann-machine formulation
of the attention and Hopfield energy. For equal-norm memories, the exact target
studied here is a Parzen kernel density estimator with bandwidth
$h=\beta^{-1/2}$ \cite{parzenEstimationProbability1962}. Kernel-density
sampling at that bandwidth is therefore the same base probability model, not a
competing generator. Nearest-neighbor interpolation, matched memory
perturbation, and fitted Gaussian mixtures remain useful controls because they
change the component allocation, covariance, or representation.

Energy-based models define a density through an energy and often estimate its
parameters using contrastive divergence
\cite{hintonTrainingProductsExperts2002}, score matching
\cite{hyvarinen2005estimation}, or noise-contrastive estimation
\cite{gutmannHyvaerinenNoiseContrastive2010}. Score-based and diffusion models
learn the score of a sequence of perturbed data distributions
\cite{songGenerativeModelingEstimating2019,songScoreBasedGenerativeModeling2021,
hoDenoisingDiffusionProbabilistic2020,sohlDicksteinDeepUnsupervisedLearning2015},
whereas normalizing flows learn an invertible transformation with a tractable
Jacobian \cite{rezendeMohamed2015}. Stochastic attention differs in that the
score of its memory energy is analytic, although this convenience also limits
the base model to a fixed memory-centered distribution. The Energy Transformer
\cite{hooverEnergyTransformer2024} embeds Hopfield-like energies in a deeper
architecture, whereas the present work isolates sampling from the tied
single-state energy. Deng et al.
\cite{dengLatentAlignment2018} introduce stochastic attention variables through
variational inference; their stochasticity and training objective are distinct
from Langevin sampling on a fixed energy. The unadjusted Langevin algorithm and
the Metropolis-adjusted Langevin algorithm are standard samplers for targets
specified through an energy or score
\cite{robertsTweedieExponentialConvergence1996,durmusNonasymptoticAnalysisUnadjusted2017}.
Our contribution is their attention-specific construction and the analysis of
the resulting memory energy. The exact mixture supplies an analytic reference
for the implementation. When a generic constraint, likelihood, or coupled
energy is added, independent ancestral sampling usually disappears but the
gradient-based construction remains available. In the base case, exact
sampling should be preferred when independent equilibrium endpoints are the
only objective.

Low-data generation also requires careful learned baselines. A variational
autoencoder with a fixed unit-Gaussian prior can place probability mass in
latent regions that were weakly constrained by a small training set
\cite{kingmaWellingAutoEncoding2014}. Geometry-aware latent models address this
problem by modeling latent support explicitly
\cite{chadebecGeometryVAE2023}, while iterative encode-decode maps can reveal
attractors and memorization in trained autoencoders
\cite{fumeroLatentDynamics2026}. The variational-autoencoder and diffusion
experiments in this paper record the tested low-data protocols, but they do not
establish general superiority over learned generators. Three subsequent studies
apply stochastic attention in this regime. The first extends training-free
generation across eight small protein-family alignments and evaluates sequence,
structural, language-model, and mutational-tolerance diagnostics
\cite{varnerTrainingFreeProtein2026}. The second uses Hopfield pattern
multiplicities to set exact mixture weights for conditional protein generation,
while separating latent conditioning from losses introduced by projection and
sequence decoding \cite{varnerConditioningProtein2026}. The third applies
multiplicity-weighted stochastic attention to a longitudinal coagulation cohort
of 23 patients, compares Langevin and direct sampling from the same
distribution, and evaluates the resulting profiles with statistical and
mechanistic tests \cite{varnerSyntheticPatient2026}. Together, these studies
use the present paper as the theoretical basis for stochastic attention while
also showing why exact base sampling and more general gradient-based sampling
should both remain part of the framework.
